\section{Exploratory Analysis Highlights}
\label{sec:eda}

Four findings from the exploratory phase shaped everything that followed; the
figures are in Appendix~\ref{app:figures} and every number is regenerated by
\code{src/report\_stats.py}.

\subsection{Growth, seasonality and drift}
\label{sec:temporal}

Monthly order volume rises roughly \textbf{eight-fold} across the window, from
about 800 orders in January 2017 to 6{,}000--7{,}500 a month through 2018,
peaking at \textbf{Black Friday, November 2017} with 7{,}544 orders, 1.5 times
the median month (Figure~\ref{fig:volume}a). Purchasing is weekday-heavy, peaks
at 16:00 and is 22\% weekend; lead time is right-skewed with a median of 10.2
days and a long tail (Figure~\ref{fig:rhythm}).

The most consequential temporal finding is not seasonality but \textbf{drift}.
Mean lead time falls from 16.9 days in February 2018 to 7.7 in August 2018 ---
Olist materially improved fulfilment over twenty months
(Figure~\ref{fig:volume}b) --- and across the chronological split the mean
falls by \textbf{5.34 days} between training and test. Under a level shift this
size, $R^2$ against a baseline fitted on the past is negative even for a useful
model, so mean absolute error leads the regression results
\citep{hyndman2006accuracy}, and periodic refitting is a deployment requirement
rather than an optimisation \citep{gama2014drift}.

\subsection{Lateness: volatile, with identifiable shocks}
\label{sec:eda-lateness}

Lateness behaves differently: it is \textbf{volatile rather than trending},
swinging from 1.4\% (June 2018) to 21.4\% (March 2018). The Phase 1 report
attributed the spikes to Black Friday and to the May 2018 truckers' strike; the
evaluators pointed out that the figure did not support the second claim, and
\code{src/lateness\_anatomy.py} re-examined the series at three resolutions
(Appendix~\ref{app:lateness}, Figure~\ref{fig:lateness}). Black Friday 2017 is
a genuine spike (14.3\%). The largest peaks, February and March 2018 (16.0\% and
21.4\%), are a network-wide \textbf{carrier-side slowdown}: transit from carrier
handoff to doorstep lengthened from 8.7 to 13.2 days while sellers' handling
time did not move, and it coincided with the platform \emph{tightening} its
promised window from 26.3 days in January to 22.0 in March. The truckers'
strike of 21 May to 1 June 2018 \citep{bloomberg2018strike,
tecnoblog2018correios} is real but invisible at monthly resolution: orders
caught in transit were 16.0\% late against a 5.0\% average over the preceding
five weeks, and the platform then lengthened its promise from 19.5 to 36.8
days, which is why June 2018 is the quietest month in the series. None of these
shocks is an attribute of an order, so nothing available at checkout can
anticipate them --- the ceiling Section~\ref{sec:bands} states in advance.

\subsection{Geography, products and payment}

The marketplace is geographically \textbf{lopsided}, and that asymmetry drives
the delivery problem: \textbf{59.7\%} of sellers sit in S\~{a}o Paulo state
against \textbf{42.0\%} of customers, only 35.7\% of orders stay within the
customer's own state, and the median customer--seller distance is
\textbf{434\,km} (Figure~\ref{fig:geography}). Distance is the strongest
checkout-time correlate of lead time ($r = 0.39$), ahead of the platform's own
estimate ($r = 0.38$). The full correlation matrix (Figure~\ref{fig:correlations})
adds three facts the models inherit: promise and distance correlate 0.52, so the
promise is largely a distance rule; freight tracks weight (0.64) rather than
price; and review score correlates $-0.33$ with lead time and essentially zero
with price or basket size --- the delivery, not the purchase, drives the rating.
Demand is long-tailed across 74 categories; the typical order is small (median
item value R\$86.90) but freight is not a rounding error (median R\$17.17, 22\%
of item value), and 51.4\% of orders are paid in instalments
(Figure~\ref{fig:products}).

\subsection{Reviews, and the relationship that dominates}
\label{sec:reviews}

Review scores are severely imbalanced --- \textbf{77.1\%} are four or five stars
--- exactly the imbalance the brief warns about; the interesting structure is
what produces the minority. A late delivery raises the probability of a 1--2
star review from \textbf{9.2\%} to \textbf{54.0\%}, a \textbf{5.9-fold} increase,
monotone in how late the parcel is (Figure~\ref{fig:reviews}). This is the
clearest relationship in the dataset, and it reproduces at order level the
logistics--rating link \citet{deshpande2023logistics} establish across 15
million orders. It also has a methodological consequence: a model predicting the
review \emph{at checkout} cannot see the variable that matters most.

\subsection{From candidates to the two problems}
\label{sec:candidates}

Three candidate problems were probed on the chronological, customer-grouped
split with untuned \code{scikit-learn} pipelines (Figure~\ref{fig:candidates};
Table~\ref{tab:checklist} in Appendix~\ref{app:dictionary}). \emph{Delivery
performance framed twice} became Problem 1: the regression reached MAE 3.89
days against 6.65 for a constant, while the direct late classifier was weak
(PR-AUC 0.121 against a 0.074 floor) and always predicting ``on time'' scored
0.926 accuracy while catching nothing. \emph{Low review at delivery} became
Problem 2, the strongest classifier found (PR-AUC 0.402 against a 0.110 floor);
the same target predicted \emph{at checkout} reached only 0.212, which made the
prediction moment a design decision rather than a detail. \emph{Repeat purchase
within 30 days} was rejected: a 1.8\% base rate, balanced accuracy 0.517, and
right-censoring that leaves no labelled orders in the test window at the more
natural 90-day horizon.
